\section{Discussion}
\label{sec:discussion}

In the CS1 holdout (139 students), the discrimination AUC of the 30-second change at $u$, in the matched comparison, fell
within the band $[0.45, 0.55]$ set in advance. Equivalence here
means that each signal, used alone as a ranking score, separated stall run-ups from matched control run-ups no better
than the band allows, on average over the pairs weighted within writers. It does not mean that typing does not change
before a stall, that stalls and controls did not both change, that a combination of signals could not separate them, or
that the signals are uninformative for every writer.

For research and practice that builds triggers for help from keystrokes, the result concerns one quantity: the change of
this estimator's two signals over the 30~s up to the last character input before a stall. It does not address what
happens later in the pause, where both signals rose (Section~\ref{sec:data}), nor a firing rule that uses the length of
the pause itself. The registered reading for a trigger is the threshold table: at each writer's 90th percentile of
controls, 10.3\% and 9.3\% of stall run-ups exceeded the threshold, against 9.6\% of controls (Section~\ref{sec:res-main}).
A conversion made after the band was set, on seen writers, under a model that shifts the stall values by a constant
and fires when a pause in the no-event stratum reaches 2~s (Appendix~\ref{app:threshold}), puts the share of firings
on pauses that become stalls at 8.2--8.9\% at an AUC of 0.50 and 8.9--10.0\% at 0.55, a rise of 0.1--1.8 points. The
low base rate sets most of this: even at 0.70 the share was 12.9--18.4\%, and a trigger that waits longer in the pause
meets a higher base rate. In that sense the 30-second change gave no basis for a trigger on this input; a trigger whose
threshold is fixed independently and whose hit rate and false alarms per hour are measured could behave differently. Those who plan to use an estimator of the same type as a trigger can
check its run-up signal against a band set in advance \cite{lakens2017equivalence} on held-out writers
\cite{gardner2018evaluating}; our data do not show the harm of omitting the check.

\subsection{How far the instrument explains the equivalence}

The equivalence could reflect the instrument rather than the typing: the input was rebuilt from edit events with a
synthesised release, the coefficients were not estimated on these data, and the levels discriminated less in CS1 (0.524
and 0.521) than in the essays (0.569 and 0.580). The positive control shows that the estimator is not inert on this
input, but also that changes of the size observed between stalls and controls barely move its 30-second changes, and
that its Stuck-state probability does not register corrections as friction; many of its scaled inputs sat at their
floor or ceiling, so a small change often leaves them unchanged. The main intervals exclude differences larger than
about 0.025 from 0.5 (a description made after the data were seen), and at this resolution the instrument cannot show
whether the run-up differs by the observed amount. The learners point to a difference of another kind: stalls were
preceded by less typing, a level rather than a change, partly because of how a stall is defined, and the estimator's
30-second changes are not built to carry it. The result is therefore a
statement about this instrument on this input, not about the absence of differences in the typing before stalls.

\subsection{What was established and what remains possible}

What the registration establishes is limited to the statistic above being inside the band in the matched comparison,
including after removing stalls that contain events inside the pause. The weak response to small changes, the floor of
the scaled inputs and the level difference are explorations; they bound the reading but were not tested. Whether unmatched evaluation makes run-up rises look large
cannot be judged here: the unmatched comparison showed no rise, and in the learner description the matched comparison
was higher than the unmatched windows (Section~\ref{sec:res-learners}). Nor can the opposite reading, that matching
removed a real run-up difference \cite{mansournia2018case}, be supported or excluded. The result is consistent with the
parallel view of planning (Section~\ref{sec:intro}) but is not a test of it, because a serial view that allows planning
to start late in the run-up would predict the same, and the view concerns language boundaries in prose rather than long
stops in code. The direction slightly below 0.5 and its candidate causes are
discussed in Appendix~\ref{app:sensitivity}.

\subsection{Calibration and reporting}

The decision rule declared equivalence more often than its nominal level in one of its two versions
(Section~\ref{sec:design}), as is known for cluster-level inference with few or unequal clusters
\cite{cameron2008bootstrap,cameron2015practitioner,mackinnon2023cluster}; designs with few effective writers may
benefit from a similar calibration. In other fields, unmatched controls have made accuracy look
larger \cite{lijmer1999empirical,pepe2008pivotal}; stating what the comparison partner was matched on lets readers judge
a run-up signal, although our data do not show an effect of the way of matching.

\subsection{Materials}

No new data were collected; the analyses used de-identified public datasets under the terms set by their authors. We
release the preregistration, the deviation, the start records and the results registration (all public on the Open
Science Framework), and an archive (version 0.1.0, \url{https://doi.org/10.5281/zenodo.24014043}) with the estimator and
the replay tool cut out at the version used, the code of the evaluation, the explorations and the conversion of the
public data, the seed and the computational environment, the manifest of the verified files, the conditions committed
before each computation (in Japanese, with an English summary) and the aggregate results. On every call that the
paper's pipelines made on the seen CS1 2019, essay and co-writing data, the released estimator returned output
byte-identical to the build used for the analysis; the calls of the 2021 data and of the positive control were not
compared, and rerunning the comparison needs that build, which is not distributed. The archive gives the steps to obtain
the public data and to rerun the confirmation on the holdout; run with the released estimator, they reproduced the
released result byte for byte. The order of the records can be checked from outside only
through the timestamps of the registrations; the commit times of the conditions come from the authors' private history,
and six files of the manifest are documents of that history, three of which (the preregistration text and its English
check) are attached to the registration. The CS1 files are released by their authors on request, and we agreed to publish
only aggregate values, so neither the data nor the per-pause estimate rows are redistributed. The list of holdout writers
will be added if the authors of the data agree; otherwise its SHA-256 and the procedure that drew it are given.
